\begin{ouverture}[courbe]{19}{La proportionnalité}{Partie 5 · Modéliser et organiser}
{Deux grandeurs · Fois plus · Fois moins}
\objectif{Reconnaître une situation de proportionnalité.}
\objectif{Raisonner avec « fois plus » et « fois moins ».}
\objectif{Résoudre des problèmes de proportionnalité en contexte.}
\prerequis{Multiplication, division et nombres décimaux.}
\end{ouverture}
\pagestyle{otMain}\essentiel\fontsize{11}{15}\selectfont
\phantomsection\label{cours-19-1}
\essSection{1}{Reconnaître une situation proportionnelle}
\begin{definition}Deux grandeurs sont en situation de proportionnalité quand, si l’une est multipliée par un nombre, l’autre est multipliée par le même nombre : 2 fois plus de baguettes, 2 fois plus cher ; 3 fois moins de farine, 3 fois moins de crêpes.\end{definition}
\begin{exemple}Le prix de cahiers identiques est proportionnel à leur nombre. En revanche, la taille d’un enfant n’est pas proportionnelle à son âge : à 10 ans, on ne mesure pas deux fois sa taille de 5 ans.\end{exemple}
\begin{attention}Un prix avec une part fixe n’est pas proportionnel : avec 10 euros d’abonnement puis 2 euros par séance, 10 séances ne coûtent pas deux fois plus que 5 séances. Vérifier le « fois plus » avant de l’utiliser.\end{attention}
\begin{savoirfaire}[Dire le raisonnement]Rédiger le raisonnement avec une phrase : « Si j’achète 3 fois plus de pains, alors je paie 3 fois plus. » La phrase dit la grandeur, le sens et le nombre de fois.\end{savoirfaire}
\clearpage
\phantomsection\label{cours-19-2}
\essSection{2}{Raisonner de proportion en proportion}
\begin{exemple}Quatre pains aux raisins coûtent 7 euros. Douze pains, c’est 3 fois plus de pains : ils coûtent 3 fois plus, soit $7\times3=21$ euros. Le prix d’un seul pain n’était pas nécessaire.\end{exemple}
\begin{savoirfaire}[Multiplier ou diviser les deux grandeurs]« Fois moins » fonctionne de même : si 10 photocopies coûtent 3 euros, 5 photocopies, c’est 2 fois moins, donc $3:2=1{,}50$ euro. Chercher le lien multiplicatif entre les deux quantités données.\end{savoirfaire}
\begin{exemple}Une fontaine régulière remplit 40 L en 5 min. En 20 min, soit 4 fois plus de temps, elle remplit $40\times4=160$ L. Le débit régulier garantit la proportionnalité.\end{exemple}
\begin{attention}Toujours contrôler la vraisemblance : plus de pains doit coûter plus cher ; moins de temps doit donner moins d’eau. Un résultat qui diminue quand la quantité augmente signale une erreur.\end{attention}

\clearpage\rubrique{Découverte}
\begin{blocLivre}{ex-19-1}\begin{exercice}[1]{Proportionnel ou pas ?}
Dire si chaque situation relève de la proportionnalité : le prix de cahiers identiques selon leur nombre ; la taille d’un enfant selon son âge.
\end{exercice}\end{blocLivre}
\begin{blocLivre}{ex-19-2}\begin{exercice}[1]{Deux fois plus}
Un kilogramme de pommes coûte $2{,}40$ euros. Combien coûtent 2 kg, puis 4 kg ?
\end{exercice}\end{blocLivre}
\begin{blocLivre}{ex-19-3}\begin{exercice}[1]{La recette multipliée}
Pour 4 personnes, une recette demande 250 g de riz. Quelle masse pour 8 personnes ? Pour 12 personnes ?
\end{exercice}\end{blocLivre}
\begin{blocLivre}{ex-19-4}\begin{exercice}[1]{Les pains aux raisins}
Quatre pains aux raisins coûtent 7 euros. Combien coûtent 12 pains aux raisins ?
\end{exercice}\end{blocLivre}
\begin{blocLivre}{ex-19-5}\begin{exercice}[1]{Les cahiers du bureau}
Cinq cahiers identiques coûtent 8 euros. Combien coûtent 15 cahiers ?
\end{exercice}\end{blocLivre}

\rubrique{Entraînement}
\begin{blocLivre}{ex-19-6}\begin{exercice}[2]{Fois moins}
Dix photocopies coûtent 3 euros. Combien coûtent 30 photocopies ? Et 5 photocopies ?
\end{exercice}\end{blocLivre}
\begin{blocLivre}{ex-19-7}\begin{exercice}[2]{Le pot de peinture}
Six litres de peinture couvrent 30 m² de mur. Quelle surface couvrent 3 litres ? Et 18 litres ?
\end{exercice}\end{blocLivre}
\begin{blocLivre}{ex-19-8}\begin{exercice}[2]{À allure régulière}
Un train roule à allure régulière : en 2 heures, il parcourt 320 km.\begin{questions}
\item Quelle distance parcourt-il en 1 heure ?
\item Quelle distance parcourt-il en 6 heures ?
\end{questions}
\end{exercice}\end{blocLivre}
\begin{blocLivre}{ex-19-9}\begin{exercice}[2]{Les crêpes du goûter}
Avec 250 g de farine, on fait 12 crêpes. Combien de crêpes avec 500 g ? Quelle farine pour 36 crêpes ?
\end{exercice}\end{blocLivre}
\begin{blocLivre}{ex-19-10}\begin{exercice}[2]{L’abonnement piégé}
La piscine coûte 10 euros d’abonnement puis 2 euros par entrée. Paul dit : « 10 entrées coûtent deux fois plus cher que 5 entrées. » Vérifier.
\end{exercice}\end{blocLivre}

\rubrique{Problèmes}
\begin{blocLivre}{ex-19-11}\begin{exercice}[3]{La fontaine du square}
Une fontaine au débit régulier remplit un arrosoir de 40 L en 5 minutes. Quel volume remplit-elle en 20 minutes ? En combien de temps remplit-elle 120 L ?
\end{exercice}\end{blocLivre}
\begin{blocLivre}{ex-19-12}\begin{exercice}[3]{Le marché}
Trois kilogrammes de poires coûtent $7{,}50$ euros. Combien coûtent 9 kg ? Et 6 kg ?
\end{exercice}\end{blocLivre}
\begin{blocLivre}{ex-19-13}\begin{exercice}[3]{Deux boulangeries}
Chez Bleuet, 4 croissants coûtent 5 euros. Chez Coquelicot, 8 croissants coûtent $9{,}60$ euros. Où 8 croissants sont-ils les moins chers ?
\end{exercice}\end{blocLivre}
\begin{blocLivre}{ex-19-14}\begin{exercice}[3]{Le ruban au mètre}
Quatre mètres de ruban coûtent 6 euros. Lisa veut 10 m : elle calcule le prix de 2 m, puis multiplie par 5. Détailler et conclure.
\end{exercice}\end{blocLivre}
\begin{blocLivre}{ex-19-15}\begin{exercice}[3]{Le carré qui grandit}
On double le côté d’un carré. Son périmètre double-t-il ? Son aire double-t-elle ? Tester avec un carré de 3 carreaux de côté.
\end{exercice}\end{blocLivre}

\clearpage
